\chapter*{ABSTRAK}
\addcontentsline{toc}{chapter}{ABSTRAK}

\begin{center}
    \textbf{Implementasi dan Integrasi Frontend Dashboard Monitoring Berbasis Web untuk Otomasi Dosis Water Treatment Plant} \\
    \vspace{0.4cm}
    oleh \\
    \vspace{0.4cm}
    Matthew Nicholas Gunawan\\
    18222058\\
\end{center}

\begin{spacing}{1.0}
Instalasi Pengolahan Air atau \textit{Water Treatment Plant} (WTP) memerlukan
pemantauan kualitas air serta proses pemberian dosis bahan kimia agar operator dapat mengenali
perubahan kondisi secara tepat. Pemantauan yang bergantung pada instrumen dan
catatan terpisah menyulitkan operator dalam mengamati parameter pH, kekeruhan,
dan \textit{Total Dissolved Solids} (TDS), serta kondisi bahan kimia dan pompa
dosis. Tugas akhir ini berfokus pada implementasi dan integrasi frontend
dashboard berbasis web untuk sistem otomasi dosis WTP ITB Jatinangor.

Frontend dikembangkan menggunakan Next.js, React, dan TypeScript serta
terintegrasi dengan backend melalui REST API. Dashboard dirancang untuk
menyediakan autentikasi operator, ringkasan kepatuhan terhadap Permenkes, visualisasi data sensor,
riwayat dan ekspor data, pemantauan alarm, status pompa, pengaturan offset dosis,
serta peringatan data kedaluwarsa. TanStack Query digunakan untuk mengelola data
server, sedangkan dukungan \textit{Progressive Web App} memungkinkan aplikasi
dipasang dan digunakan melalui perangkat desktop maupun perangkat seluler. Ruang lingkup
tugas akhir dibatasi pada implementasi frontend, sedangkan perangkat sensor, backend, basis data, dan
algoritma kendali dosis berada di luar ruang lingkup.

Evaluasi dilakukan melalui pengujian fungsional terhadap kebutuhan dan alur
utama frontend. Hasil pengujian menunjukkan bahwa seluruh 20 kasus uji fungsional
berhasil dipenuhi. Hasil tersebut memverifikasi bahwa frontend mampu menangani autentikasi, 
menampilkan ringkasan pemantauan dan analitik sensor, mengekspor riwayat data, 
serta memvalidasi interaksi alarm dan offset dosis sesuai kebutuhan sistem.

\noindent\textbf{Kata kunci:} dashboard, frontend, instalasi pengolahan air, Next.js, \textit{Progressive Web App}.
\end{spacing}
